\section{The six-point conic criterion}
\label{sec:logical}

Section~\ref{sec:exact-logical} already determines the fixed-coordinate
transversal logical group at every length.  In the six-coordinate case, the
same coding-theoretic test has a classical geometric interpretation: the
conductor \(\Cond(C,C^\perp)\) is nonzero exactly when the six projective
points lie on a conic.

Throughout this section \(q\) is an odd prime.  Regard any one coordinate of
an \(\AME(6,q)\) stabilizer tensor as an input and the other five as outputs.
This is a \([[5,1,3]]_q\) encoding isometry.  In prime local dimension, a
one-qudit Clifford acts symplectically through \(\mathrm{SL}_2(q)\), so the
fixed-coordinate symplectic stabilizer of the label space is
\[
 \mathcal K_C=\{(F_1,\ldots,F_6)\in\mathrm{SL}_2(q)^6:
                 (\bigoplus_iF_i)L_C=L_C\}.
\]
For an encoder view, the \emph{fixed-coordinate logical symplectic image} is
the image of \(\mathcal K_C\) on its input block.

Theorem~\ref{thm:diagonal-isodual-transversal-group} expresses the dichotomy
as diagonal equivalence to the dual.  In projective geometry, the same
condition is called fixed-label Gale self-association.  The labeled Gale
transform consists of the six columns of any generator matrix for \(\ker H\).
Fixed-label self-association means that this configuration is projectively
equivalent to the columns of \(H\) without permuting their labels.  The next
lemma identifies this condition with the conic locus.

\begin{lemma}[Fixed-label self-association]
\label{lem:six-arc-self-association}
Let \(H\) be a rank-three \(3\times6\) matrix whose columns are a six-arc,
and put \(C=\ker H\).  There is a nonsingular diagonal matrix \(S\) with
\(SC=C^\perp\) if and only if the six columns of \(H\) lie on a conic.
\end{lemma}

\begin{proof}
Let \(G\) be a \(3\times6\) generator matrix for \(C\).  Its columns are
the labeled Gale transform of the columns of \(H\).  Since
\(C^\perp=\operatorname{row}(H)\),
\[
 SC=C^\perp
 \quad\Longleftrightarrow\quad
 \operatorname{row}(GS)=\operatorname{row}(H)
 \quad\Longleftrightarrow\quad
 AGS=H\ \text{ for some }A\in\mathrm{GL}_3(q).
\]
Thus the condition is precisely fixed-label projective equivalence to the
Gale transform.  This fixed locus can be identified directly over the finite
field.  Put
\[
 V_i=(x_i^2,y_i^2,z_i^2,x_iy_i,x_iz_i,y_iz_i)
\]
and let \(V\) be the \(6\)-by-\(6\) matrix with rows \(V_i\).  The six
entries are the degree-two monomials in three variables, equivalently the
independent coefficients of a symmetric \(3\)-by-\(3\) matrix.  Thus a
vector in the kernel of \(V\) is a quadratic form vanishing on all six
points, and \(\det V=0\) is exactly the conic condition.

Any five arc points impose independent conditions on conics.  Indeed, a
relation among at most five rows of \(V\) would give
\(H_I D H_I^T=0\) for a nonsingular diagonal \(D\) on its support \(I\).
For \(|I|\geq3\), this would put the three-dimensional row space of
\(H_I D\) inside \(\ker H_I\), which has dimension at most two; smaller
supports are immediate.  Hence any relation
\(\sum_i d_i V_i=0\) has every \(d_i\ne0\).

Writing \(D=\diag(d_i)\), that relation is \(HDH^T=0\).  Therefore
\(\operatorname{row}(HD)\subseteq\ker H=C\), and equality follows from
dimension three.  Thus \(CD^{-1}=C^\perp\).  Conversely, a nonsingular
diagonal equivalence \(SC=C^\perp\) gives
\(HS^{-1}H^T=0\), hence a relation among the rows of \(V\) and a conic
through the six columns.  The conic is irreducible because the columns form
an arc.  This is the finite-field fixed-label form of the classical
six-point self-association theorem
\cite[Section~2.3]{HowardMillsonSnowdenVakil2008}.
\end{proof}

\begin{theorem}[Six-point transversal-group criterion]
\label{thm:logical-phase}
For an equal-phase CSS state from a six-arc over an odd prime field, the
fixed-coordinate logical symplectic image is
\begin{equation}\label{eq:six-point-linear-logical-group}
\begin{cases}
\mathrm{SL}_2(q),&\text{if the arc is conic/GRS},\\
T=\{\diag(a,a^{-1}):a\in\F_q^\times\},&\text{if the arc is nonconic}.
\end{cases}
\end{equation}
Thus every encoder view of a GRS tensor realizes the full one-qudit logical
symplectic group from fixed-coordinate symmetries.  Off the GRS locus, the
image is the split torus.
\end{theorem}

\begin{proof}
Theorem~\ref{thm:diagonal-isodual-transversal-group} gives
\(\mathrm{SL}_2(q)\) exactly when \(C\) is diagonally equivalent to its dual,
and \(T\) otherwise.  Lemma~\ref{lem:six-arc-self-association} identifies
this condition with the conic locus.
\end{proof}

Adding logical Paulis recovers the exact fixed-coordinate projective logical
groups in Theorem~\ref{thm:diagonal-isodual-transversal-group}.  This is an
operational distinction, not a classical automorphism count: the non-GRS
group has order \(q^2(q-1)\), whereas the GRS group has order
\(q^2\lvert\mathrm{SL}_2(q)\rvert\).

Allowing coordinate permutations is a separate problem.  A permutation that
exchanges the CSS \(X\)- and \(Z\)-spaces can supply the nontrivial Weyl-group
element normalizing \(T\), but acts on a fixed encoder view only when it fixes
the input coordinate; Appendix~\ref{sec:party-extensions} treats this
extension.  Appendix~\ref{sec:six-point-example} gives an optional worked
example separating syndrome symmetry from the parity-check conic criterion.
